% =========================================================
% Slides Beamer Metropolis 16:9 -- Chapter 5 (ARMA models)
% Time Series Analysis -- Aymen Ben Brik (aymen.benbrik@esprit.tn)
%
% Migrated from _archives/source-slides/Stationnarite.tex
% sections "Modeles AR/MA/ARMA" (1083-1575), "Analyse des residus"
% (1575-1762), "Forecasting" (1762-1950).
% =========================================================
\documentclass[aspectratio=169]{beamer}

% =========================================================
% Beamer preamble — Time Series Analysis module
% Theme : Metropolis, aspect ratio 16:9
% Author : Aymen Ben Brik (aymen.benbrik@esprit.tn)
% Esprit School of Business
%
% Include from each chapterX/slides.tex via :
%   % =========================================================
% Beamer preamble — Time Series Analysis module
% Theme : Metropolis, aspect ratio 16:9
% Author : Aymen Ben Brik (aymen.benbrik@esprit.tn)
% Esprit School of Business
%
% Include from each chapterX/slides.tex via :
%   % =========================================================
% Beamer preamble — Time Series Analysis module
% Theme : Metropolis, aspect ratio 16:9
% Author : Aymen Ben Brik (aymen.benbrik@esprit.tn)
% Esprit School of Business
%
% Include from each chapterX/slides.tex via :
%   \input{../preamble/slides-preamble.tex}
% AFTER the line \documentclass[aspectratio=169]{beamer}.
% =========================================================

\usepackage[utf8]{inputenc}
\usepackage[T1]{fontenc}
\usepackage[english]{babel}
\usepackage{lmodern}
\usepackage{amsmath, amssymb, mathtools, bm}
\usepackage{graphicx}
\usepackage{booktabs}
\usepackage{array}
\usepackage{multirow}
\usepackage{colortbl}
\usepackage{xcolor}
\usepackage{tikz}
\usetikzlibrary{positioning, arrows.meta, calc, shapes, fit, backgrounds, matrix, decorations.pathreplacing}
\usepackage{listings}
\usepackage[most]{tcolorbox}

% --- Theme ---
\usetheme{metropolis}
\setbeamertemplate{navigation symbols}{}
\metroset{block=fill}

% --- Esprit colour palette ---
\definecolor{espritBlue}{RGB}{30, 60, 110}
\definecolor{espritAccent}{RGB}{200, 80, 30}
\setbeamercolor{progress bar}{fg=espritAccent}
\setbeamercolor{title separator}{fg=espritAccent}
\setbeamercolor{frametitle}{bg=espritBlue}

% --- Code listings (Python + R, with UTF-8 literate) ---
\definecolor{codebg}{RGB}{248,248,248}
\definecolor{codeframe}{RGB}{220,220,220}
\definecolor{keyword}{RGB}{0,82,155}
\definecolor{comment}{RGB}{88,110,117}
\definecolor{string}{RGB}{133,153,0}

\lstdefinestyle{pythonstyle}{
  language=Python,
  basicstyle=\ttfamily\scriptsize,
  keywordstyle=\color{keyword}\bfseries,
  commentstyle=\color{comment}\itshape,
  stringstyle=\color{string},
  backgroundcolor=\color{codebg},
  frame=single,
  rulecolor=\color{codeframe},
  frameround=tttt,
  breaklines=true,
  showstringspaces=false,
  tabsize=2,
  literate=
    {é}{{\'e}}1 {è}{{\`e}}1 {ê}{{\^e}}1 {à}{{\`a}}1 {â}{{\^a}}1
    {ç}{{\c{c}}}1 {²}{{$^2$}}1 {°}{{${}^\circ$}}1
    {α}{{$\alpha$}}1 {β}{{$\beta$}}1 {γ}{{$\gamma$}}1
    {δ}{{$\delta$}}1 {θ}{{$\theta$}}1 {λ}{{$\lambda$}}1
    {μ}{{$\mu$}}1 {σ}{{$\sigma$}}1 {ε}{{$\varepsilon$}}1
    {ρ}{{$\rho$}}1 {τ}{{$\tau$}}1 {φ}{{$\phi$}}1
    {≤}{{$\leq$}}1 {≥}{{$\geq$}}1 {≈}{{$\approx$}}1
}
\lstdefinestyle{Rstyle}{
  language=R,
  basicstyle=\ttfamily\scriptsize,
  keywordstyle=\color{keyword}\bfseries,
  commentstyle=\color{comment}\itshape,
  stringstyle=\color{string},
  backgroundcolor=\color{codebg},
  frame=single,
  rulecolor=\color{codeframe},
  frameround=tttt,
  breaklines=true,
  showstringspaces=false,
  tabsize=2
}
\lstset{style=pythonstyle}

% --- Common macros ---
\input{../preamble/macros.tex}

% --- tcolorbox boxes for slides ---
\definecolor{boxTh}{RGB}{120, 30, 30}
\definecolor{boxProp}{RGB}{30, 100, 60}
\definecolor{boxDef}{RGB}{30, 60, 110}
\definecolor{boxEx}{RGB}{180, 120, 0}

\newtcolorbox{infobox}[1][]{enhanced, colback=espritBlue!5, colframe=espritBlue,
  fonttitle=\bfseries\small,
  title={\ifstrempty{#1}{}{#1}},
  boxrule=0.4pt, arc=2pt, left=4pt, right=4pt, top=2pt, bottom=2pt}
\newtcolorbox{warnbox}[1][]{enhanced, colback=espritAccent!5, colframe=espritAccent,
  fonttitle=\bfseries\small,
  title={\ifstrempty{#1}{}{#1}},
  boxrule=0.4pt, arc=2pt, left=4pt, right=4pt, top=2pt, bottom=2pt}
\newtcolorbox{thbox}[1][]{enhanced, colback=boxTh!5, colframe=boxTh,
  fonttitle=\bfseries\small,
  title={Theorem\ifstrempty{#1}{}{ -- #1}},
  boxrule=0.5pt, arc=2pt, left=4pt, right=4pt, top=2pt, bottom=2pt}
\newtcolorbox{propbox}[1][]{enhanced, colback=boxProp!5, colframe=boxProp,
  fonttitle=\bfseries\small,
  title={Proposition\ifstrempty{#1}{}{ -- #1}},
  boxrule=0.5pt, arc=2pt, left=4pt, right=4pt, top=2pt, bottom=2pt}
\newtcolorbox{defbox}[1][]{enhanced, colback=boxDef!5, colframe=boxDef,
  fonttitle=\bfseries\small,
  title={Definition\ifstrempty{#1}{}{ -- #1}},
  boxrule=0.5pt, arc=2pt, left=4pt, right=4pt, top=2pt, bottom=2pt}
\newtcolorbox{exbox}[1][]{enhanced, colback=boxEx!5, colframe=boxEx,
  fonttitle=\bfseries\small,
  title={Example\ifstrempty{#1}{}{ -- #1}},
  boxrule=0.5pt, arc=2pt, left=4pt, right=4pt, top=2pt, bottom=2pt}

% AFTER the line \documentclass[aspectratio=169]{beamer}.
% =========================================================

\usepackage[utf8]{inputenc}
\usepackage[T1]{fontenc}
\usepackage[english]{babel}
\usepackage{lmodern}
\usepackage{amsmath, amssymb, mathtools, bm}
\usepackage{graphicx}
\usepackage{booktabs}
\usepackage{array}
\usepackage{multirow}
\usepackage{colortbl}
\usepackage{xcolor}
\usepackage{tikz}
\usetikzlibrary{positioning, arrows.meta, calc, shapes, fit, backgrounds, matrix, decorations.pathreplacing}
\usepackage{listings}
\usepackage[most]{tcolorbox}

% --- Theme ---
\usetheme{metropolis}
\setbeamertemplate{navigation symbols}{}
\metroset{block=fill}

% --- Esprit colour palette ---
\definecolor{espritBlue}{RGB}{30, 60, 110}
\definecolor{espritAccent}{RGB}{200, 80, 30}
\setbeamercolor{progress bar}{fg=espritAccent}
\setbeamercolor{title separator}{fg=espritAccent}
\setbeamercolor{frametitle}{bg=espritBlue}

% --- Code listings (Python + R, with UTF-8 literate) ---
\definecolor{codebg}{RGB}{248,248,248}
\definecolor{codeframe}{RGB}{220,220,220}
\definecolor{keyword}{RGB}{0,82,155}
\definecolor{comment}{RGB}{88,110,117}
\definecolor{string}{RGB}{133,153,0}

\lstdefinestyle{pythonstyle}{
  language=Python,
  basicstyle=\ttfamily\scriptsize,
  keywordstyle=\color{keyword}\bfseries,
  commentstyle=\color{comment}\itshape,
  stringstyle=\color{string},
  backgroundcolor=\color{codebg},
  frame=single,
  rulecolor=\color{codeframe},
  frameround=tttt,
  breaklines=true,
  showstringspaces=false,
  tabsize=2,
  literate=
    {é}{{\'e}}1 {è}{{\`e}}1 {ê}{{\^e}}1 {à}{{\`a}}1 {â}{{\^a}}1
    {ç}{{\c{c}}}1 {²}{{$^2$}}1 {°}{{${}^\circ$}}1
    {α}{{$\alpha$}}1 {β}{{$\beta$}}1 {γ}{{$\gamma$}}1
    {δ}{{$\delta$}}1 {θ}{{$\theta$}}1 {λ}{{$\lambda$}}1
    {μ}{{$\mu$}}1 {σ}{{$\sigma$}}1 {ε}{{$\varepsilon$}}1
    {ρ}{{$\rho$}}1 {τ}{{$\tau$}}1 {φ}{{$\phi$}}1
    {≤}{{$\leq$}}1 {≥}{{$\geq$}}1 {≈}{{$\approx$}}1
}
\lstdefinestyle{Rstyle}{
  language=R,
  basicstyle=\ttfamily\scriptsize,
  keywordstyle=\color{keyword}\bfseries,
  commentstyle=\color{comment}\itshape,
  stringstyle=\color{string},
  backgroundcolor=\color{codebg},
  frame=single,
  rulecolor=\color{codeframe},
  frameround=tttt,
  breaklines=true,
  showstringspaces=false,
  tabsize=2
}
\lstset{style=pythonstyle}

% --- Common macros ---
% =========================================================
% Common math macros — Time Series Analysis module
% Author : Aymen Ben Brik (aymen.benbrik@esprit.tn)
% Esprit School of Business
% =========================================================

% --- Standard sets ---
\providecommand{\R}{\mathbb{R}}
\providecommand{\N}{\mathbb{N}}
\providecommand{\Z}{\mathbb{Z}}
\providecommand{\Q}{\mathbb{Q}}
\providecommand{\C}{\mathbb{C}}

% --- Probability / statistics ---
\providecommand{\E}{\mathbb{E}}
\providecommand{\Prob}{\mathbb{P}}
\providecommand{\Var}{\mathrm{Var}}
\providecommand{\Cov}{\mathrm{Cov}}
\providecommand{\Corr}{\mathrm{Corr}}
\providecommand{\indep}{\perp\!\!\!\perp}
\providecommand{\iid}{\overset{\text{i.i.d.}}{\sim}}
\providecommand{\given}{\,\big\vert\,}

% --- Estimators / hat ---
\providecommand{\hatm}{\widehat{m}}
\providecommand{\hatsig}{\widehat{\sigma}}
\providecommand{\hatrho}{\widehat{\rho}}
\providecommand{\hatphi}{\widehat{\phi}}
\providecommand{\hattheta}{\widehat{\theta}}

% --- Time series notation ---
\providecommand{\Lop}{L}                      % lag operator
\providecommand{\Diff}{\Delta}                % difference operator
\providecommand{\AR}[1]{\mathrm{AR}(#1)}
\providecommand{\MAm}[1]{\mathrm{MA}(#1)}
\providecommand{\ARMA}[2]{\mathrm{ARMA}(#1,#2)}
\providecommand{\ARIMA}[3]{\mathrm{ARIMA}(#1,#2,#3)}
\providecommand{\SARIMA}[6]{\mathrm{SARIMA}(#1,#2,#3)(#4,#5,#6)}
\providecommand{\ARCH}[1]{\mathrm{ARCH}(#1)}
\providecommand{\GARCH}[2]{\mathrm{GARCH}(#1,#2)}
\providecommand{\WN}{\mathrm{WN}}              % white noise
\providecommand{\MSE}{\mathrm{MSE}}
\providecommand{\MAE}{\mathrm{MAE}}
\providecommand{\RMSE}{\mathrm{RMSE}}
\providecommand{\AIC}{\mathrm{AIC}}
\providecommand{\BIC}{\mathrm{BIC}}
\providecommand{\ACF}{\mathrm{ACF}}
\providecommand{\PACF}{\mathrm{PACF}}

% --- Linear algebra ---
\providecommand{\transp}[1]{#1^{\!\top}}
\providecommand{\norm}[1]{\left\lVert #1 \right\rVert}
\providecommand{\abs}[1]{\left\lvert #1 \right\rvert}
\providecommand{\inner}[2]{\langle #1, #2 \rangle}
\providecommand{\diag}{\mathrm{diag}}
\providecommand{\tr}{\mathrm{tr}}

% --- Bold vectors / matrices ---
\providecommand{\vx}{\bm{x}}
\providecommand{\vy}{\bm{y}}
\providecommand{\vz}{\bm{z}}
\providecommand{\vh}{\bm{h}}
\providecommand{\vu}{\bm{u}}
\providecommand{\vv}{\bm{v}}
\providecommand{\vw}{\bm{w}}
\providecommand{\vb}{\bm{b}}
\providecommand{\vW}{\bm{W}}
\providecommand{\vSigma}{\bm{\Sigma}}
\providecommand{\veps}{\bm{\varepsilon}}

% --- Author identity (reused on titlepages) ---
\providecommand{\auteurNom}{Aymen Ben Brik}
\providecommand{\auteurEmail}{aymen.benbrik@esprit.tn}
\providecommand{\auteurInstitution}{Esprit School of Business}
\providecommand{\auteurRole}{Module head}
\providecommand{\moduleNom}{Time Series Analysis}
\providecommand{\anneeUniversitaire}{2025--2026}


% --- tcolorbox boxes for slides ---
\definecolor{boxTh}{RGB}{120, 30, 30}
\definecolor{boxProp}{RGB}{30, 100, 60}
\definecolor{boxDef}{RGB}{30, 60, 110}
\definecolor{boxEx}{RGB}{180, 120, 0}

\newtcolorbox{infobox}[1][]{enhanced, colback=espritBlue!5, colframe=espritBlue,
  fonttitle=\bfseries\small,
  title={\ifstrempty{#1}{}{#1}},
  boxrule=0.4pt, arc=2pt, left=4pt, right=4pt, top=2pt, bottom=2pt}
\newtcolorbox{warnbox}[1][]{enhanced, colback=espritAccent!5, colframe=espritAccent,
  fonttitle=\bfseries\small,
  title={\ifstrempty{#1}{}{#1}},
  boxrule=0.4pt, arc=2pt, left=4pt, right=4pt, top=2pt, bottom=2pt}
\newtcolorbox{thbox}[1][]{enhanced, colback=boxTh!5, colframe=boxTh,
  fonttitle=\bfseries\small,
  title={Theorem\ifstrempty{#1}{}{ -- #1}},
  boxrule=0.5pt, arc=2pt, left=4pt, right=4pt, top=2pt, bottom=2pt}
\newtcolorbox{propbox}[1][]{enhanced, colback=boxProp!5, colframe=boxProp,
  fonttitle=\bfseries\small,
  title={Proposition\ifstrempty{#1}{}{ -- #1}},
  boxrule=0.5pt, arc=2pt, left=4pt, right=4pt, top=2pt, bottom=2pt}
\newtcolorbox{defbox}[1][]{enhanced, colback=boxDef!5, colframe=boxDef,
  fonttitle=\bfseries\small,
  title={Definition\ifstrempty{#1}{}{ -- #1}},
  boxrule=0.5pt, arc=2pt, left=4pt, right=4pt, top=2pt, bottom=2pt}
\newtcolorbox{exbox}[1][]{enhanced, colback=boxEx!5, colframe=boxEx,
  fonttitle=\bfseries\small,
  title={Example\ifstrempty{#1}{}{ -- #1}},
  boxrule=0.5pt, arc=2pt, left=4pt, right=4pt, top=2pt, bottom=2pt}

% AFTER the line \documentclass[aspectratio=169]{beamer}.
% =========================================================

\usepackage[utf8]{inputenc}
\usepackage[T1]{fontenc}
\usepackage[english]{babel}
\usepackage{lmodern}
\usepackage{amsmath, amssymb, mathtools, bm}
\usepackage{graphicx}
\usepackage{booktabs}
\usepackage{array}
\usepackage{multirow}
\usepackage{colortbl}
\usepackage{xcolor}
\usepackage{tikz}
\usetikzlibrary{positioning, arrows.meta, calc, shapes, fit, backgrounds, matrix, decorations.pathreplacing}
\usepackage{listings}
\usepackage[most]{tcolorbox}

% --- Theme ---
\usetheme{metropolis}
\setbeamertemplate{navigation symbols}{}
\metroset{block=fill}

% --- Esprit colour palette ---
\definecolor{espritBlue}{RGB}{30, 60, 110}
\definecolor{espritAccent}{RGB}{200, 80, 30}
\setbeamercolor{progress bar}{fg=espritAccent}
\setbeamercolor{title separator}{fg=espritAccent}
\setbeamercolor{frametitle}{bg=espritBlue}

% --- Code listings (Python + R, with UTF-8 literate) ---
\definecolor{codebg}{RGB}{248,248,248}
\definecolor{codeframe}{RGB}{220,220,220}
\definecolor{keyword}{RGB}{0,82,155}
\definecolor{comment}{RGB}{88,110,117}
\definecolor{string}{RGB}{133,153,0}

\lstdefinestyle{pythonstyle}{
  language=Python,
  basicstyle=\ttfamily\scriptsize,
  keywordstyle=\color{keyword}\bfseries,
  commentstyle=\color{comment}\itshape,
  stringstyle=\color{string},
  backgroundcolor=\color{codebg},
  frame=single,
  rulecolor=\color{codeframe},
  frameround=tttt,
  breaklines=true,
  showstringspaces=false,
  tabsize=2,
  literate=
    {é}{{\'e}}1 {è}{{\`e}}1 {ê}{{\^e}}1 {à}{{\`a}}1 {â}{{\^a}}1
    {ç}{{\c{c}}}1 {²}{{$^2$}}1 {°}{{${}^\circ$}}1
    {α}{{$\alpha$}}1 {β}{{$\beta$}}1 {γ}{{$\gamma$}}1
    {δ}{{$\delta$}}1 {θ}{{$\theta$}}1 {λ}{{$\lambda$}}1
    {μ}{{$\mu$}}1 {σ}{{$\sigma$}}1 {ε}{{$\varepsilon$}}1
    {ρ}{{$\rho$}}1 {τ}{{$\tau$}}1 {φ}{{$\phi$}}1
    {≤}{{$\leq$}}1 {≥}{{$\geq$}}1 {≈}{{$\approx$}}1
}
\lstdefinestyle{Rstyle}{
  language=R,
  basicstyle=\ttfamily\scriptsize,
  keywordstyle=\color{keyword}\bfseries,
  commentstyle=\color{comment}\itshape,
  stringstyle=\color{string},
  backgroundcolor=\color{codebg},
  frame=single,
  rulecolor=\color{codeframe},
  frameround=tttt,
  breaklines=true,
  showstringspaces=false,
  tabsize=2
}
\lstset{style=pythonstyle}

% --- Common macros ---
% =========================================================
% Common math macros — Time Series Analysis module
% Author : Aymen Ben Brik (aymen.benbrik@esprit.tn)
% Esprit School of Business
% =========================================================

% --- Standard sets ---
\providecommand{\R}{\mathbb{R}}
\providecommand{\N}{\mathbb{N}}
\providecommand{\Z}{\mathbb{Z}}
\providecommand{\Q}{\mathbb{Q}}
\providecommand{\C}{\mathbb{C}}

% --- Probability / statistics ---
\providecommand{\E}{\mathbb{E}}
\providecommand{\Prob}{\mathbb{P}}
\providecommand{\Var}{\mathrm{Var}}
\providecommand{\Cov}{\mathrm{Cov}}
\providecommand{\Corr}{\mathrm{Corr}}
\providecommand{\indep}{\perp\!\!\!\perp}
\providecommand{\iid}{\overset{\text{i.i.d.}}{\sim}}
\providecommand{\given}{\,\big\vert\,}

% --- Estimators / hat ---
\providecommand{\hatm}{\widehat{m}}
\providecommand{\hatsig}{\widehat{\sigma}}
\providecommand{\hatrho}{\widehat{\rho}}
\providecommand{\hatphi}{\widehat{\phi}}
\providecommand{\hattheta}{\widehat{\theta}}

% --- Time series notation ---
\providecommand{\Lop}{L}                      % lag operator
\providecommand{\Diff}{\Delta}                % difference operator
\providecommand{\AR}[1]{\mathrm{AR}(#1)}
\providecommand{\MAm}[1]{\mathrm{MA}(#1)}
\providecommand{\ARMA}[2]{\mathrm{ARMA}(#1,#2)}
\providecommand{\ARIMA}[3]{\mathrm{ARIMA}(#1,#2,#3)}
\providecommand{\SARIMA}[6]{\mathrm{SARIMA}(#1,#2,#3)(#4,#5,#6)}
\providecommand{\ARCH}[1]{\mathrm{ARCH}(#1)}
\providecommand{\GARCH}[2]{\mathrm{GARCH}(#1,#2)}
\providecommand{\WN}{\mathrm{WN}}              % white noise
\providecommand{\MSE}{\mathrm{MSE}}
\providecommand{\MAE}{\mathrm{MAE}}
\providecommand{\RMSE}{\mathrm{RMSE}}
\providecommand{\AIC}{\mathrm{AIC}}
\providecommand{\BIC}{\mathrm{BIC}}
\providecommand{\ACF}{\mathrm{ACF}}
\providecommand{\PACF}{\mathrm{PACF}}

% --- Linear algebra ---
\providecommand{\transp}[1]{#1^{\!\top}}
\providecommand{\norm}[1]{\left\lVert #1 \right\rVert}
\providecommand{\abs}[1]{\left\lvert #1 \right\rvert}
\providecommand{\inner}[2]{\langle #1, #2 \rangle}
\providecommand{\diag}{\mathrm{diag}}
\providecommand{\tr}{\mathrm{tr}}

% --- Bold vectors / matrices ---
\providecommand{\vx}{\bm{x}}
\providecommand{\vy}{\bm{y}}
\providecommand{\vz}{\bm{z}}
\providecommand{\vh}{\bm{h}}
\providecommand{\vu}{\bm{u}}
\providecommand{\vv}{\bm{v}}
\providecommand{\vw}{\bm{w}}
\providecommand{\vb}{\bm{b}}
\providecommand{\vW}{\bm{W}}
\providecommand{\vSigma}{\bm{\Sigma}}
\providecommand{\veps}{\bm{\varepsilon}}

% --- Author identity (reused on titlepages) ---
\providecommand{\auteurNom}{Aymen Ben Brik}
\providecommand{\auteurEmail}{aymen.benbrik@esprit.tn}
\providecommand{\auteurInstitution}{Esprit School of Business}
\providecommand{\auteurRole}{Module head}
\providecommand{\moduleNom}{Time Series Analysis}
\providecommand{\anneeUniversitaire}{2025--2026}


% --- tcolorbox boxes for slides ---
\definecolor{boxTh}{RGB}{120, 30, 30}
\definecolor{boxProp}{RGB}{30, 100, 60}
\definecolor{boxDef}{RGB}{30, 60, 110}
\definecolor{boxEx}{RGB}{180, 120, 0}

\newtcolorbox{infobox}[1][]{enhanced, colback=espritBlue!5, colframe=espritBlue,
  fonttitle=\bfseries\small,
  title={\ifstrempty{#1}{}{#1}},
  boxrule=0.4pt, arc=2pt, left=4pt, right=4pt, top=2pt, bottom=2pt}
\newtcolorbox{warnbox}[1][]{enhanced, colback=espritAccent!5, colframe=espritAccent,
  fonttitle=\bfseries\small,
  title={\ifstrempty{#1}{}{#1}},
  boxrule=0.4pt, arc=2pt, left=4pt, right=4pt, top=2pt, bottom=2pt}
\newtcolorbox{thbox}[1][]{enhanced, colback=boxTh!5, colframe=boxTh,
  fonttitle=\bfseries\small,
  title={Theorem\ifstrempty{#1}{}{ -- #1}},
  boxrule=0.5pt, arc=2pt, left=4pt, right=4pt, top=2pt, bottom=2pt}
\newtcolorbox{propbox}[1][]{enhanced, colback=boxProp!5, colframe=boxProp,
  fonttitle=\bfseries\small,
  title={Proposition\ifstrempty{#1}{}{ -- #1}},
  boxrule=0.5pt, arc=2pt, left=4pt, right=4pt, top=2pt, bottom=2pt}
\newtcolorbox{defbox}[1][]{enhanced, colback=boxDef!5, colframe=boxDef,
  fonttitle=\bfseries\small,
  title={Definition\ifstrempty{#1}{}{ -- #1}},
  boxrule=0.5pt, arc=2pt, left=4pt, right=4pt, top=2pt, bottom=2pt}
\newtcolorbox{exbox}[1][]{enhanced, colback=boxEx!5, colframe=boxEx,
  fonttitle=\bfseries\small,
  title={Example\ifstrempty{#1}{}{ -- #1}},
  boxrule=0.5pt, arc=2pt, left=4pt, right=4pt, top=2pt, bottom=2pt}


\title[\moduleNom\ -- Ch.5]{ARMA Models}
\subtitle{AR $\cdot$ MA $\cdot$ ARMA $\cdot$ Estimation $\cdot$ Diagnostics $\cdot$ Forecasting}
\author{\auteurNom \\ \footnotesize\auteurEmail}
\institute{\auteurInstitution}
\date{\anneeUniversitaire}

\begin{document}

\begin{frame}\titlepage\end{frame}

\begin{frame}{Learning objectives}
  By the end of this chapter, you will be able to:
  \begin{itemize}
    \item write an AR$(p)$, MA$(q)$, ARMA$(p,q)$ model in lag-operator
          form;
    \item state the stationarity and invertibility conditions in
          terms of polynomial roots;
    \item read the ACF/PACF \emph{signatures} of AR, MA, ARMA;
    \item estimate an ARMA model by maximum likelihood;
    \item validate the model with residual diagnostics
          (Ljung--Box, normality);
    \item produce $h$-step-ahead forecasts with confidence
          intervals.
  \end{itemize}
\end{frame}

\begin{frame}{Chapter outline}
  \begin{enumerate}
    \item Autoregressive process AR$(p)$
    \item Moving-average process MA$(q)$
    \item ARMA$(p,q)$
    \item Box--Jenkins identification (ACF / PACF signatures)
    \item Estimation: MLE, conditional sum of squares
    \item Model selection: AIC / BIC
    \item Residual diagnostics
    \item Forecasting: point forecast and confidence interval
  \end{enumerate}
\end{frame}

% =====================================================================
\section{Autoregressive process AR$(p)$}

\begin{frame}{1.~Autoregressive process AR$(p)$}
  \begin{block}{Definition}
    An autoregressive process of order $p$ satisfies
    \[
      X_t = \phi_1 X_{t-1} + \phi_2 X_{t-2} + \dots + \phi_p X_{t-p}
            + \varepsilon_t,
      \qquad \varepsilon_t \sim WN(0, \sigma^2).
    \]
    The current value depends on a finite history of \emph{past
    values} of the series, plus an innovation $\varepsilon_t$.
  \end{block}
  \pause
  Special case AR$(1)$:
  \[
    X_t = \phi X_{t-1} + \varepsilon_t.
  \]
  \begin{itemize}
    \item $\phi$ controls persistence.
    \item Large $|\phi|$ $\Rightarrow$ slow mean reversion.
  \end{itemize}
\end{frame}

\begin{frame}{AR$(p)$ in lag-operator form}
  Recall $L X_t = X_{t-1}$ and define the AR polynomial
  \[
    \Phi(z) = 1 - \phi_1 z - \phi_2 z^2 - \dots - \phi_p z^p.
  \]
  The AR$(p)$ writes compactly as
  \[
    \Phi(L)\, X_t = \varepsilon_t.
  \]
  \begin{block}{Stationarity condition}
    AR$(p)$ admits a (unique, causal) stationary solution if and only
    if all roots of $\Phi(z) = 0$ lie \emph{strictly outside} the
    unit disc, i.e.\ $|z_i| > 1$ for $i = 1, \dots, p$.
  \end{block}
  AR$(1)$: root is $z = 1/\phi$, condition becomes $|\phi| < 1$.
\end{frame}

\begin{frame}{AR$(1)$: moments and ACF}
  Assume $|\phi| < 1$. Iterating
  $X_t = \phi X_{t-1} + \varepsilon_t$ backward gives
  \[
    X_t = \sum_{j=0}^{\infty} \phi^j \varepsilon_{t-j}
    \qquad\text{(MA$(\infty)$ representation)}.
  \]
  \pause
  Hence
  \[
    \E[X_t] = 0,
    \qquad
    \Var(X_t) = \frac{\sigma^2}{1 - \phi^2},
    \qquad
    \rho(h) = \phi^{|h|}.
  \]
  \begin{itemize}
    \item ACF \textbf{decays geometrically} with rate $\phi$.
    \item PACF \textbf{cuts off} after lag~$1$:
          $\alpha(1) = \phi$, $\alpha(h) = 0$ for $h \geq 2$.
  \end{itemize}
\end{frame}

\begin{frame}{AR$(2)$: stationarity triangle}
  \[
    X_t = \phi_1 X_{t-1} + \phi_2 X_{t-2} + \varepsilon_t.
  \]
  Stationarity $\Leftrightarrow$ $(\phi_1, \phi_2)$ in the triangle
  \[
    \phi_1 + \phi_2 < 1, \qquad
    \phi_2 - \phi_1 < 1, \qquad
    -1 < \phi_2 < 1.
  \]
  \pause
  Yule--Walker equations give the ACF recursively:
  \[
    \rho(h) = \phi_1 \rho(h-1) + \phi_2 \rho(h-2),
    \qquad h \geq 1.
  \]
  with $\rho(0) = 1$,
  $\rho(1) = \phi_1 / (1 - \phi_2)$.
\end{frame}

% =====================================================================
\section{Moving-average process MA$(q)$}

\begin{frame}{2.~Moving-average process MA$(q)$}
  \begin{block}{Definition}
    \[
      X_t = \varepsilon_t + \theta_1 \varepsilon_{t-1}
            + \dots + \theta_q \varepsilon_{t-q},
      \qquad \varepsilon_t \sim WN(0, \sigma^2).
    \]
    The current value is a finite linear combination of \emph{present
    and past shocks}.
  \end{block}
  \pause
  Compact form with $\Theta(z) = 1 + \theta_1 z + \dots + \theta_q z^q$:
  \[
    X_t = \Theta(L)\, \varepsilon_t.
  \]
  \alert{Every MA$(q)$ is stationary} (finite linear combination of
  white noise).
\end{frame}

\begin{frame}{MA$(q)$: moments and ACF}
  \[
    \E[X_t] = 0,
    \qquad
    \Var(X_t) = \sigma^2 \bigl(1 + \theta_1^2 + \dots + \theta_q^2\bigr).
  \]
  Autocovariance:
  \[
    \gamma(h) =
    \begin{cases}
      \sigma^2 \bigl(\theta_h + \theta_1 \theta_{h+1} + \dots
                     + \theta_{q-h}\theta_q\bigr) & 1 \leq h \leq q, \\[2pt]
      0                                          & h > q.
    \end{cases}
  \]
  \pause
  \begin{itemize}
    \item ACF \textbf{cuts off} at lag $q$ -- the defining MA
          signature.
    \item PACF decays geometrically (no cutoff).
  \end{itemize}
\end{frame}

\begin{frame}{Invertibility of MA$(q)$}
  An MA$(q)$ is \emph{invertible} if it can be rewritten as an
  AR$(\infty)$:
  \[
    \varepsilon_t = X_t - \pi_1 X_{t-1} - \pi_2 X_{t-2} - \dots
  \]
  with $\sum |\pi_j| < \infty$.
  \begin{block}{Invertibility condition}
    All roots of $\Theta(z) = 0$ lie \emph{strictly outside} the unit
    disc: $|z_j| > 1$ for $j = 1, \dots, q$.
  \end{block}
  \pause
  MA$(1)$ example: $X_t = \varepsilon_t + \theta \varepsilon_{t-1}$
  is invertible iff $|\theta| < 1$. Without invertibility, two
  parameter sets give the same ACF -- the model is not identifiable.
\end{frame}

% =====================================================================
\section{ARMA$(p,q)$}

\begin{frame}{3.~ARMA$(p,q)$}
  \begin{block}{Definition}
    \[
      \Phi(L)\, X_t = \Theta(L)\, \varepsilon_t,
    \]
    or, expanded,
    \[
      X_t = \sum_{i=1}^{p} \phi_i X_{t-i}
          + \varepsilon_t + \sum_{j=1}^{q} \theta_j \varepsilon_{t-j}.
    \]
  \end{block}
  \pause
  \begin{itemize}
    \item AR part captures persistence in past values.
    \item MA part captures persistence of past shocks.
    \item Special cases: ARMA$(p, 0) = $ AR$(p)$,
          ARMA$(0, q) = $ MA$(q)$.
  \end{itemize}
\end{frame}

\begin{frame}{ARMA$(p,q)$: existence and identifiability}
  \begin{block}{Conditions}
    \begin{itemize}
      \item \textbf{Stationarity}: roots of $\Phi(z) = 0$ outside the
            unit disc.
      \item \textbf{Invertibility}: roots of $\Theta(z) = 0$ outside
            the unit disc.
      \item \textbf{Identifiability}: $\Phi$ and $\Theta$ have no common
            root (no AR/MA cancellation).
    \end{itemize}
  \end{block}
  \pause
  Under these three conditions, the ARMA process is a stationary
  causal invertible solution of the equation, with absolutely
  summable MA$(\infty)$ representation $X_t = \Phi^{-1}(L)\Theta(L)
  \varepsilon_t$.
\end{frame}

% =====================================================================
\section{Box--Jenkins identification}

\begin{frame}{4.~ACF / PACF signatures}
  \begin{center}
  \renewcommand{\arraystretch}{1.4}
  \begin{tabular}{lll}
  \toprule
  Model       & ACF $\rho(h)$         & PACF $\alpha(h)$ \\
  \midrule
  AR$(p)$     & exponential decay     & \textbf{cuts off} after lag $p$ \\
  MA$(q)$     & \textbf{cuts off} after lag $q$ & exponential decay \\
  ARMA$(p,q)$ & exponential decay (after lag $q$) & exponential decay (after lag $p$) \\
  \bottomrule
  \end{tabular}
  \end{center}
  \vspace{0.2cm}
  \pause
  \begin{block}{Box--Jenkins identification rule}
    \begin{itemize}
      \item PACF cuts off at lag $p$ \& ACF decays $\Rightarrow$ AR$(p)$.
      \item ACF cuts off at lag $q$ \& PACF decays $\Rightarrow$ MA$(q)$.
      \item Both decay $\Rightarrow$ ARMA -- compare AIC/BIC across orders.
    \end{itemize}
  \end{block}
\end{frame}

\begin{frame}{Cutoff vs decay: visual cue}
  Bartlett bands $\pm 1.96/\sqrt{T}$ delimit ``zero''.
  \begin{itemize}
    \item ``Cuts off after lag $k$'' means
          $\widehat\rho(h)$ (or $\widehat\alpha(h)$) drops inside the
          band for all $h > k$.
    \item ``Decays exponentially'' means $\widehat\rho(h)$ stays
          outside the band for many lags but shrinks.
  \end{itemize}
  \pause
  \alert{Caveat:} on real data, the picture is often noisy.
  Identification gives \emph{candidate} orders -- the final choice
  uses information criteria and residual diagnostics.
\end{frame}

% =====================================================================
\section{Estimation}

\begin{frame}{5.~Estimation -- Yule--Walker for AR$(p)$}
  Multiplying $X_t = \sum \phi_i X_{t-i} + \varepsilon_t$ by
  $X_{t-h}$ and taking expectations gives the Yule--Walker system:
  \[
    \rho(h) = \phi_1 \rho(h-1) + \dots + \phi_p \rho(h-p),
    \qquad h = 1, \dots, p.
  \]
  In matrix form $\bm R \bm\phi = \bm r$:
  \[
    \begin{pmatrix}
      1 & \rho(1) & \cdots & \rho(p-1) \\
      \rho(1) & 1 & \cdots & \rho(p-2) \\
      \vdots & & \ddots & \vdots \\
      \rho(p-1) & \cdots & \rho(1) & 1
    \end{pmatrix}
    \begin{pmatrix}\phi_1 \\ \vdots \\ \phi_p\end{pmatrix}
    =
    \begin{pmatrix}\rho(1) \\ \vdots \\ \rho(p)\end{pmatrix}.
  \]
  Plug $\widehat\rho$ in place of $\rho$ to obtain
  $\widehat\phi_{\text{YW}}$. Method-of-moments style; consistent but
  not efficient.
\end{frame}

\begin{frame}{Estimation -- Conditional sum of squares}
  Conditioning on the first $p$ observations and assuming
  $\varepsilon_t = 0$ for $t \leq 0$:
  \[
    \widehat{\bm\theta}_{\text{CSS}}
    = \arg\min_{\bm\theta}
      \sum_{t = p+1}^{T} \widehat\varepsilon_t(\bm\theta)^2,
  \]
  with $\widehat\varepsilon_t$ obtained recursively from the model
  equation. Equivalent to MLE under Gaussian innovations and large
  $T$, easier to compute.
\end{frame}

\begin{frame}{Estimation -- Maximum likelihood}
  Under $\varepsilon_t \iid \mathcal N(0, \sigma^2)$, the joint
  likelihood of $(X_1, \dots, X_T)$ is multivariate normal with
  covariance $\Gamma_T(\bm\theta)$:
  \[
    \log \mathcal L(\bm\theta, \sigma^2)
    = -\tfrac{T}{2}\log(2\pi)
      -\tfrac{1}{2}\log|\Gamma_T|
      -\tfrac{1}{2} \bm X^\top \Gamma_T^{-1} \bm X.
  \]
  \pause
  Maximised numerically (state-space + Kalman filter; the standard
  backend of \texttt{statsmodels.SARIMAX}). Provides asymptotic
  standard errors and $t$-statistics for each coefficient.
\end{frame}

\begin{frame}[fragile]{Estimation -- Python (statsmodels)}
\begin{lstlisting}[language=Python]
from statsmodels.tsa.arima.model import ARIMA

model = ARIMA(y, order=(2, 0, 1))   # ARMA(2,1)
fit   = model.fit()
print(fit.summary())
print("AIC =", fit.aic, " BIC =", fit.bic)
print("Coefs:", fit.params)
\end{lstlisting}
  \begin{itemize}
    \item \texttt{order=(p, d, q)} -- here $d = 0$ for ARMA.
    \item \texttt{fit.summary()} prints estimates, standard errors,
          $z$-statistics, $p$-values, AIC, BIC, Ljung--Box.
    \item \texttt{fit.resid} gives the residuals; use them in
          diagnostics.
  \end{itemize}
\end{frame}

% =====================================================================
\section{Model selection}

\begin{frame}{6.~Information criteria}
  Compare candidate orders by penalised log-likelihood:
  \[
    \mathrm{AIC} = -2\, \log \mathcal L + 2\,k,
    \qquad
    \mathrm{BIC} = -2\, \log \mathcal L + k \log T,
  \]
  where $k = p + q + 1$ is the number of free parameters.
  \pause
  \begin{itemize}
    \item AIC favours predictive accuracy; can over-fit on small $T$.
    \item BIC penalises complexity more (consistent under the true
          order assumption).
  \end{itemize}
  \alert{Practical rule:} when AIC and BIC disagree, prefer the
  parsimonious BIC choice and check it with residual diagnostics.
\end{frame}

% =====================================================================
\section{Residual diagnostics}

\begin{frame}{7.~Residual diagnostics}
  Once $\widehat{\bm\theta}$ is obtained, compute the residuals
  $\widehat\varepsilon_t = X_t - \widehat X_{t \mid t-1}$.
  A well-fitted model satisfies:
  \begin{itemize}
    \item $\widehat\varepsilon_t$ should be \textbf{white noise};
    \item residual ACF inside the Bartlett band;
    \item Ljung--Box test does not reject $H_0$ (white noise);
    \item residuals approximately normal (QQ-plot, Jarque--Bera).
  \end{itemize}
  \pause
  Failure of any test $\Rightarrow$ revisit identification: try a
  larger AR/MA order, or look for a structural break / volatility
  effect (Chapter 6).
\end{frame}

\begin{frame}{Ljung--Box test on residuals}
  Test
  \[
    Q_{\text{LB}}(m)
    = T(T+2) \sum_{h=1}^{m}
      \frac{\widehat\rho_{\widehat\varepsilon}(h)^2}{T - h}.
  \]
  Under $H_0$ (residuals are WN), and after fitting an ARMA$(p,q)$,
  \[
    Q_{\text{LB}}(m) \xrightarrow{d} \chi^2_{m - p - q}
    \qquad\text{(degrees of freedom adjusted for fitted parameters)}.
  \]
  \pause
  \alert{Convention.} Test at $m = 12, 24$ for monthly data; $m =
  10, 20$ for short series. Reject if $p$-value $< 0.05$.
\end{frame}

% =====================================================================
\section{Forecasting}

\begin{frame}{8.~Forecasting an ARMA$(p,q)$}
  Best linear forecast at horizon $h$:
  \[
    \widehat X_{T+h \mid T}
    = \E\bigl[X_{T+h} \mid X_T, X_{T-1}, \dots\bigr].
  \]
  For AR$(1)$:
  $\widehat X_{T+h \mid T} = \phi^h X_T \to 0$ as $h \to \infty$.
  \pause
  Forecast error variance grows with $h$:
  \[
    \Var\bigl(X_{T+h} - \widehat X_{T+h \mid T}\bigr)
    = \sigma^2 \sum_{j=0}^{h-1} \psi_j^2,
  \]
  where $\psi_j$ are the MA$(\infty)$ coefficients
  ($X_t = \sum \psi_j \varepsilon_{t-j}$).
\end{frame}

\begin{frame}{Forecast confidence intervals}
  Under Gaussian innovations, a $95\%$ confidence interval for
  $X_{T+h}$ is
  \[
    \widehat X_{T+h \mid T}
    \;\pm\; 1.96\, \sigma\,
        \sqrt{\sum_{j=0}^{h-1} \psi_j^2}.
  \]
  \pause
  \begin{itemize}
    \item For stationary ARMA, the CI width tends to a finite limit
          (standard deviation of the unconditional distribution).
    \item For ARIMA (with $d \geq 1$), the CI grows without bound --
          consistent with the unit root.
  \end{itemize}
\end{frame}

\begin{frame}[fragile]{Forecasting -- Python (statsmodels)}
\begin{lstlisting}[language=Python]
fit = ARIMA(y, order=(2, 0, 1)).fit()

fc      = fit.get_forecast(steps=12)
mean    = fc.predicted_mean
ci_low  = fc.conf_int(alpha=0.05).iloc[:, 0]
ci_high = fc.conf_int(alpha=0.05).iloc[:, 1]

import matplotlib.pyplot as plt
plt.plot(y, label="observed")
plt.plot(range(len(y), len(y) + 12), mean, "r-", label="forecast")
plt.fill_between(range(len(y), len(y) + 12), ci_low, ci_high,
                 color="red", alpha=0.2)
plt.legend(); plt.show()
\end{lstlisting}
\end{frame}

% =====================================================================
\section{Box--Jenkins workflow}

\begin{frame}{The Box--Jenkins workflow}
  \begin{enumerate}
    \item \textbf{Stationarise} (Chapter 4): differencing,
          log-transform if needed.
    \item \textbf{Identify} candidate $(p, q)$ from ACF / PACF
          signatures.
    \item \textbf{Estimate} each candidate by MLE.
    \item \textbf{Select} the model with lowest AIC / BIC, subject to
          parsimony.
    \item \textbf{Diagnose} the residuals (Ljung--Box, QQ-plot).
    \item \textbf{Forecast} with confidence bands.
  \end{enumerate}
  \pause
  \alert{If diagnostics fail}, return to step~2 with a richer order or
  consider seasonal effects (SARIMA) or volatility (Chapter 6).
\end{frame}

\begin{frame}{Take-home messages}
  \begin{itemize}
    \item ARMA combines past values (AR) and past shocks (MA).
    \item Stationarity = AR roots outside the unit disc;
          invertibility = MA roots outside the unit disc.
    \item ACF / PACF tell you the order: \emph{cutoff} reveals AR or
          MA; \emph{double decay} suggests ARMA.
    \item Estimation: MLE in \texttt{statsmodels}; selection by AIC /
          BIC; residuals must be white noise.
    \item Forecast intervals widen with horizon at a rate dictated
          by the MA$(\infty)$ representation.
  \end{itemize}
\end{frame}

\end{document}
